\section{Программная реализация прототипа системы обнаружения и ИИ-злоумышленника}
\label{ch:impl}

\noindent
\textit{Содержание главы соответствует третьему этапу постановки
задачи (см.~ВКР, Этап~3): разработка программного прототипа системы
обнаружения аномалий (проектирование архитектуры, реализация модулей
предобработки и формирования окон, ядра детекции, минимального
интерфейса), обучение и оптимизация моделей с подготовкой реальных и
синтетических данных, разработка ИИ-злоумышленника и библиотеки
атакующих сценариев.}\\

\subsection{Стек и общая структура программного проекта}

Прототип реализован на платформе Windows 11 в среде VS Code в виде
Python-пакета \texttt{diploma\_nids}. Технологический стек:

\begin{itemize}
    \item Python 3.10+ (тестировался на 3.13.11);
    \item PyTorch 2.x (CPU-сборка) --- глубокое обучение;
    \item scikit-learn 1.x --- классические baseline-модели;
    \item XGBoost --- gradient boosting baseline;
    \item Pydantic v2 --- валидация конфигов и схем данных;
    \item PyYAML --- параметризация всех стадий pipeline;
    \item FastAPI + uvicorn --- REST API инференса;
    \item Streamlit --- визуальная панель мониторинга;
    \item pytest --- юнит-тестирование (9 модулей);
    \item Make --- оркестрация сборки и воспроизводимости.
\end{itemize}

Общая структура проекта приведена в \tabref{tab:project-layout}.
Общий объём исходного кода --- около \textbf{4230 строк Python}.
В колонке ``Каталог'' префикс \texttt{src/diploma\_nids/} для подпакетов
\texttt{data/}, \texttt{models/}, \texttt{training/} и т.\,д. опущен ради
компактности (физический путь --- \texttt{src/diploma\_nids/<каталог>}).

{
\renewcommand{\arraystretch}{1.25}
\setlength{\tabcolsep}{4pt}
\small
\begin{longtable}{@{}p{2.6cm}p{4.6cm}p{7.5cm}@{}}
\caption{Структура программного проекта \texttt{diploma\_nids}}
\label{tab:project-layout}\\
\toprule
\textbf{Каталог} & \textbf{Назначение} & \textbf{Ключевые модули} \\
\midrule
\endfirsthead
\multicolumn{3}{@{}l@{}}{\textit{Продолжение таблицы~\ref{tab:project-layout}}}\\
\toprule
\textbf{Каталог} & \textbf{Назначение} & \textbf{Ключевые модули} \\
\midrule
\endhead
\bottomrule
\endfoot
\texttt{configs/}      & Параметризация всех стадий (26 YAML-файлов) &
\texttt{data}, \texttt{preprocess}, \texttt{models}, \texttt{train},
\texttt{eval}, \texttt{attacker}, \texttt{pipeline} \\
\addlinespace
\texttt{data/} & Схемы, загрузка, препроцессинг, окна, splits &
\texttt{schema.py}, \texttt{load.py}, \texttt{preprocess.py},
\texttt{windowing.py}, \texttt{splits.py} \\
\addlinespace
\texttt{models/} & Модели-кандидаты &
\texttt{base.py} (registry), \texttt{cnn\_lstm.py}, \texttt{rnn\_family.py},
\texttt{cnn1d.py}, \texttt{tcn.py}, \texttt{transformer.py},
\texttt{autoencoder.py}, \texttt{classical.py} \\
\addlinespace
\texttt{training/} & Обучение DL-моделей &
\texttt{trainer.py}, \texttt{losses.py} (BCE, focal) \\
\addlinespace
\texttt{eval/} & Метрики, калибровка, drift &
\texttt{metrics.py}, \texttt{thresholding.py}, \texttt{drift.py},
\texttt{error\_analysis.py} \\
\addlinespace
\texttt{attacker/} & ИИ-злоумышленник &
\texttt{templates/} (7 шаблонов + normal), \texttt{agent.py} (FSM),
\texttt{drift\_injector.py}, \texttt{runtime.py} \\
\addlinespace
\texttt{inference/} & Online-инференс и алертинг &
\texttt{stream.py}, \texttt{alerting.py}, \texttt{service.py} (FastAPI) \\
\addlinespace
\texttt{ui/} & Streamlit-дэшборд &
\texttt{streamlit\_app.py} \\
\addlinespace
\texttt{utils/} & Воспроизводимость и IO &
\texttt{seed.py}, \texttt{io.py}, \texttt{logging.py} \\
\addlinespace
\texttt{scripts/} & Точки входа CLI & 12 нумерованных скриптов:
загрузка~UNSW, препроцессинг, обучение, оценка, прогон агента,
drift-тест, demo-стенд, генерация графиков и таблиц \\
\addlinespace
\texttt{tests/} & Юнит-тесты (9 модулей) &
\texttt{test\_utils}, \texttt{test\_configs},
\texttt{test\_preprocess}, \texttt{test\_windowing},
\texttt{test\_splits}, \texttt{test\_models\_smoke},
\texttt{test\_metrics}, \texttt{test\_attacker},
\texttt{test\_alerting} \\
\end{longtable}
}

\subsection{Подсистема обработки данных}

\subsubsection{Схема и загрузка}

Схема UNSW-NB15 формализована Pydantic-моделями (модуль
\texttt{data/schema.py}). Загрузчик \texttt{data/load.py} валидирует
наличие признаков и меток, проверяет существование официального
train/test-разделения и возвращает \texttt{DataFrame}. Аналогично
реализована загрузка CICIDS2017 для cross-evaluation.

\subsubsection{Препроцессор}

Класс \texttt{Preprocessor} реализует sklearn-совместимый интерфейс
\texttt{fit / transform / fit\_transform}. На стадии \texttt{fit}
(на обучающей выборке) формируются:

\begin{itemize}
    \item квантильные границы для обрезки выбросов
    (\(0{,}001\) и \(0{,}999\));
    \item статистики robust-scaler (медиана, IQR);
    \item one-hot категории для признаков с малой кардинальностью и
    frequency-encoding для большой кардинальности;
    \item фиксированный порядок выходных столбцов
    \texttt{output\_columns}, обеспечивающий идентичность
    представления в train/val/test/inference.
\end{itemize}

Препроцессор сериализуется в JSON (\texttt{save}) и восстанавливается
(\texttt{load}) для online-инференса.

\subsubsection{Скользящие окна}

Модуль \texttt{data/windowing.py} реализует \texttt{build\_windows},
который по векторизованному массиву признаков и меткам возвращает
тензор \((m, W, d)\) и метки окон. Параметры:
\(W = 32\), \(S = 8\); стратегия \emph{last}-агрегации меток
обоснована на этапе~\ref{ch:design} структурой UNSW-NB15.

\subsection{Реализованные модели}

В \tabref{tab:models-impl} сведены 15 реализованных моделей.
Регистрация выполняется при импорте подпакета \texttt{models}
через декоратор \texttt{@register("name")}.

\begin{table}[H]
\centering
\caption{Реализованные модели в \texttt{diploma\_nids/models/}}
\label{tab:models-impl}
\renewcommand{\arraystretch}{1.2}
\small
\begin{tabularx}{\linewidth}{@{}lllX@{}}
\toprule
\textbf{Имя} & \textbf{Тип} & \textbf{Окна} & \textbf{Файл} \\
\midrule
mlp                  & DL supervised        & нет & \texttt{models/mlp.py} \\
cnn1d                & DL supervised        & да  & \texttt{models/cnn1d.py} \\
\textbf{cnn\_lstm}   & DL supervised \emph{(proposed)} & да & \texttt{models/cnn\_lstm.py} \\
lstm                 & DL supervised        & да  & \texttt{models/rnn\_family.py} \\
gru                  & DL supervised        & да  & \texttt{models/rnn\_family.py} \\
bilstm               & DL supervised + attn & да  & \texttt{models/rnn\_family.py} \\
tcn                  & DL supervised        & да  & \texttt{models/tcn.py} \\
transformer          & DL supervised        & да  & \texttt{models/transformer.py} \\
autoencoder          & DL semi-supervised   & нет & \texttt{models/autoencoder.py} \\
vae                  & DL semi-supervised   & нет & \texttt{models/autoencoder.py} \\
logistic\_regression & classical            & нет & \texttt{models/classical.py} \\
random\_forest       & classical            & нет & \texttt{models/classical.py} \\
xgboost              & classical            & нет & \texttt{models/classical.py} \\
isolation\_forest    & classical (unsup.)   & нет & \texttt{models/classical.py} \\
ocsvm                & classical (one-class)& нет & \texttt{models/classical.py} \\
\bottomrule
\end{tabularx}
\end{table}

Все DL-модели реализованы на PyTorch и единообразно возвращают объект
\texttt{ModelOutput} с полями \texttt{logits}, \texttt{probs} и
\texttt{extras}. Архитектура CNN-LSTM:

\begin{equation*}
\begin{aligned}
\text{Input}(B, W, F)\ &\to\ \underbrace{[\text{Conv1D}\!\to\!\text{BN}\!\to\!\text{ReLU}\!\to\!\text{MaxPool}\!\to\!\text{Dropout}]\!\times\!2}_{\text{два свёрточных блока}}\\
&\to\ \text{LSTM}\ \to\ \text{MLP-head}\ \to\ \text{Sigmoid}.
\end{aligned}
\end{equation*}

\subsection{Подсистема обучения и оценки}

\textbf{Loss-функции.} В \texttt{training/losses.py} реализованы
binary cross-entropy с логитами и фокальная функция
потерь~\cite{lin2017focal} с параметрами \(\alpha\), \(\gamma\),
задаваемыми в \texttt{configs/train/default.yaml}.

\textbf{Тренер.} \texttt{training/trainer.py} реализует AdamW + cosine
LR-scheduler, gradient clipping, early stopping по
\texttt{val\_pr\_auc}, сохранение лучшего чекпойнта и журнала
истории обучения (loss, val\_pr\_auc, val\_f1). Полный набор
гиперпараметров обучения, использованный в \emph{full}-режиме при
проведении эксперимента~E1, приведён в \tabref{tab:train-hyperparams}.

\begin{table}[H]
\centering
\caption{Гиперпараметры обучения DL-моделей (конфигурация
\texttt{configs/train/full.yaml})}
\label{tab:train-hyperparams}
\renewcommand{\arraystretch}{1.2}
\begin{tabularx}{\linewidth}{@{}lXX@{}}
\toprule
\textbf{Параметр} & \textbf{Значение} & \textbf{Назначение} \\
\midrule
Оптимизатор          & AdamW            & адаптивный градиентный метод с decoupled weight decay \\
Learning rate        & $1{,}0 \cdot 10^{-3}$    & базовая скорость обучения \\
Weight decay         & $1{,}0 \cdot 10^{-4}$    & регуляризация $\ell_2$ \\
$\beta_1, \beta_2$   & $0{,}9$, $0{,}999$ & параметры Adam \\
Batch size           & 256              & размер минибатча \\
Эпохи (макс.)        & 30               & лимит обучения \\
Patience             & 7                & окно early stopping \\
Метрика мониторинга  & \texttt{val\_pr\_auc} & целевая метрика остановки (\eng{max}-режим) \\
LR-scheduler         & cosine annealing & снижение lr к $1{,}0 \cdot 10^{-6}$ \\
Gradient clipping    & $\|g\|_2 \leq 1{,}0$ & защита от взрыва градиентов \\
Loss-функция         & focal loss       & противодействие дисбалансу классов \\
$\alpha$ (focal)     & $0{,}25$         & вес положительного класса \\
$\gamma$ (focal)     & $2{,}0$          & фокусировка на «трудных» примерах \\
Class weights        & auto             & автоматический пересчёт по обучающей выборке \\
Seed                 & 42 (123, 2024)   & фиксация воспроизводимости (3 запуска для proposed) \\
Определённость       & deterministic    & PyTorch deterministic mode \\
\bottomrule
\end{tabularx}
\end{table}

\textbf{Метрики.} \texttt{eval/metrics.py} реализует Precision,
Recall, F1, FPR, MCC, ROC-AUC, PR-AUC, ECE, FP-per-time, TTD, поиск
порога по целевому FPR и по F1-max, bootstrap-CI.

\textbf{Калибровка.} \texttt{eval/thresholding.py} реализует
температурное масштабирование~\eqref{eq:temp-scaling} через L-BFGS на
валидационных логитах.

\textbf{Drift-метрики.} \texttt{eval/drift.py} реализует PSI, KL и MMD
с автобандвидтом по медианной эвристике; функция
\texttt{drift\_report} возвращает агрегат с явным индикатором
\texttt{drift\_alarm} при превышении PSI-порога \(0{,}25\).

\subsection{ИИ-злоумышленник}

\subsubsection{Шаблоны атак}

Реализованы семь параметризованных шаблонов
(\tabref{tab:attack-templates}), каждый формирует
последовательность объектов \texttt{FlowRecord} с явной привязкой к
классу \emph{attack\_cat} UNSW-NB15. Дополнительно реализован шаблон
\texttt{normal} для генерации фонового трафика.

\begin{table}[H]
\centering
\caption{Реализованная библиотека шаблонов атак}
\label{tab:attack-templates}
\renewcommand{\arraystretch}{1.2}
\setlength{\tabcolsep}{4pt}
\small
\begin{tabularx}{\linewidth}{@{}lXX@{}}
\toprule
\textbf{Шаблон} & \textbf{Параметризация} & \textbf{Соответствие классу UNSW} \\
\midrule
\texttt{ddos\_synflood} &
интенсивность \(\lambda_0\), длительность \(\Delta\), доля SYN-флагов,
распределение исходных IP &
DoS \\
\addlinespace
\texttt{port\_scan} &
скорость подключений \(\rho\), диапазон портов, режим, длительность &
Reconnaissance \\
\addlinespace
\texttt{brute\_force} &
интервал между попытками, целевой порт, доля успешных авторизаций &
Reconnaissance / Exploits \\
\addlinespace
\texttt{exploit\_payload} &
размер payload, нестандартные сочетания TCP-флагов, доля retransmissions &
Exploits \\
\addlinespace
\texttt{data\_exfil} &
длительность сессии, асимметрия \(b_{out}/b_{in}\), стабильность скорости &
Backdoors \\
\addlinespace
\texttt{internal\_recon} &
фан-аут потоков, периодичность, размер сессии &
Reconnaissance \\
\addlinespace
\texttt{worm\_lateral} &
число направленных потоков, периодичность, размер сессии &
Worms \\
\bottomrule
\end{tabularx}
\end{table}

\subsubsection{Конечно-автоматный планировщик}

\texttt{attacker/agent.py} реализует \texttt{FSMAgent} с поддержкой
11 состояний (\emph{NORMAL}, 7 \emph{ATTACK} и 3 \emph{DRIFT}),
вероятностями переходов из \texttt{configs/attacker/policy.yaml} и
минимальными длительностями состояний. Полный перечень состояний
с привязкой к шаблонам атак и классам \emph{attack\_cat} UNSW-NB15
приведён в \tabref{tab:fsm-states}. Каждый тик автомата возвращает
структуру \texttt{AgentTick} с метаданными ground truth
(\(\langle\)состояние FSM, тип атаки, тип дрейфа\(\rangle\)),
необходимыми для расчёта метрик в Эксперименте~E4.

\begin{table}[H]
\centering
\caption{Состояния конечно-автоматного планировщика
ИИ-злоумышленника}
\label{tab:fsm-states}
\renewcommand{\arraystretch}{1.2}
\setlength{\tabcolsep}{4pt}
\small
\begin{tabularx}{\linewidth}{@{}lllX@{}}
\toprule
\textbf{Состояние} & \textbf{Тип} & \textbf{Шаблон} &
\textbf{Класс UNSW-NB15} \\
\midrule
\texttt{NORMAL}         & нормальный  & \texttt{normal}         & Normal \\
\addlinespace
\texttt{ATTACK\_DDOS}   & атака       & \texttt{ddos\_synflood} & DoS \\
\texttt{ATTACK\_SCAN}   & атака       & \texttt{port\_scan}     & Reconnaissance \\
\texttt{ATTACK\_BRUTE}  & атака       & \texttt{brute\_force}   & Reconnaissance \\
\texttt{ATTACK\_EXPLOIT}& атака       & \texttt{exploit\_payload} & Exploits \\
\texttt{ATTACK\_EXFIL}  & атака       & \texttt{data\_exfil}    & Backdoors \\
\texttt{ATTACK\_INTREC} & атака       & \texttt{internal\_recon}& Reconnaissance \\
\texttt{ATTACK\_WORM}   & атака       & \texttt{worm\_lateral}  & Worms \\
\addlinespace
\texttt{DRIFT\_COV}     & дрейф       & ---                     & сдвиг масштаба/шума числовых признаков \\
\texttt{DRIFT\_PRIOR}   & дрейф       & ---                     & изменение долей классов в потоке \\
\texttt{DRIFT\_CONCEPT} & дрейф       & ---                     & маскирование сигнатурных признаков атак \\
\bottomrule
\end{tabularx}
\end{table}

\subsubsection{Drift-инжектор и runtime}

\texttt{attacker/drift\_injector.py} реализует три типа дрейфа:
\emph{covariate} (сдвиг масштаба и шума), \emph{prior} (изменение
долей классов через политику FSM), \emph{concept} (маскирование
сигнатурных признаков атаки).

\texttt{attacker/runtime.py} объединяет агент и drift-инжектор и
предоставляет два режима работы: offline (метод
\texttt{collect\_dataframe}) для эксперимента E4 и online
(\texttt{stream\_with\_callback}) для realtime-демо.

\subsection{Подсистема online-инференса и интерфейс}

\subsubsection{FastAPI-сервис}

Модуль \texttt{inference/service.py} реализует REST API:

\begin{itemize}
    \item \texttt{GET /health} --- проверка живости;
    \item \texttt{GET /info} --- сведения о загруженной модели;
    \item \texttt{POST /score} --- скоринг батча flow-записей;
    \item \texttt{POST /score-window} --- прямой скоринг готового
    окна \((W, F)\);
    \item \texttt{GET /alerts/recent} --- последние \(n\) алертов.
\end{itemize}

Загрузка артефактов происходит через переменные окружения
\texttt{DIPLOMA\_MODEL\_YAML}, \texttt{DIPLOMA\_CHECKPOINT},
\texttt{DIPLOMA\_PREPROCESSOR}, \texttt{DIPLOMA\_THRESHOLD},
что обеспечивает декларативность конфигурации.

\subsubsection{Формирование алертов}

\texttt{AlertFormer} реализует пятиступенчатую шкалу severity
(\emph{info, low, medium, high, critical}), привязанную к
нормированному превышению скоринга над порогом, и временное окно
дедупликации (\(\Delta t = 5\,\text{с}\)) для подавления шумовых
повторов.

\subsubsection{Streamlit-дэшборд}

\texttt{ui/streamlit\_app.py} реализует обновляемую с интервалом
1--2~с панель: метрики (число алертов, доля \emph{critical/high},
средний скоринг), таблицу алертов и временной ряд скорингов.
Сообщается с FastAPI-сервисом через \texttt{DIPLOMA\_API\_URL}.

\subsection{Воспроизводимость и журналирование}

\begin{itemize}
    \item Seed фиксируется через \texttt{utils/seed.py}
    (\texttt{numpy}, \texttt{random}, \texttt{torch}, \texttt{torch.cuda},
    \texttt{PYTHONHASHSEED}, \texttt{torch.use\_deterministic\_algorithms}).
    \item Артефакты обучения сохраняются в JSON-формате в
    \texttt{experiments/runs/}; чекпойнты --- в \texttt{models/}.
    \item Конфигурации экспериментов фиксируются YAML-файлами;
    \texttt{make reproduce} запускает полный цикл
    \emph{data $\to$ preprocess $\to$ train $\to$ evaluate}
    одной командой.
    \item Препроцессор сохраняется JSON-ом и применяется идентично в
    обучении и инференсе, исключая test-time leakage.
\end{itemize}

\subsection{Покрытие тестами}

Реализованы 9 модулей юнит-тестов (pytest), фиксирующих критические
инварианты программного прототипа. Состав модулей и проверяемые
свойства приведены в \tabref{tab:tests}.

\begin{table}[H]
\centering
\caption{Состав юнит-тестов программного прототипа}
\label{tab:tests}
\renewcommand{\arraystretch}{1.2}
\small
\begin{tabularx}{\linewidth}{@{}lX@{}}
\toprule
\textbf{Файл теста} & \textbf{Проверяемые свойства} \\
\midrule
\texttt{test\_utils.py}        & seed-воспроизводимость; YAML/JSON roundtrip; создание родительских директорий \\
\texttt{test\_configs.py}      & все 26 YAML-конфигов загружаются; FSM содержит 11 состояний; вероятности переходов суммируются в~1{,}0 \\
\texttt{test\_preprocess.py}   & shape выходов; отсутствие NaN/Inf; неизменность статистик при повторном fit; save/load roundtrip; обработка неизвестной категории; обрезка выбросов \\
\texttt{test\_windowing.py}    & корректные shapes; три стратегии агрегации меток; защита от коротких последовательностей; валидация входов \\
\texttt{test\_splits.py}       & сохранение временного порядка; воспроизводимость random; непересечение split-ов \\
\texttt{test\_models\_smoke.py}& 10 DL и 5 классических моделей загружаются из YAML и выполняют forward; save/load CNN-LSTM; \emph{proposed = true} \\
\texttt{test\_metrics.py}      & Precision, Recall, F1, MCC; threshold для целевого FPR; ECE снижается калибровкой; PSI/KL/MMD на одинаковых и сдвинутых распределениях; bootstrap CI \\
\texttt{test\_attacker.py}     & генерация всех 7 шаблонов; FSM эмиссия тиков; drift-инжектор сдвигает признаки; runtime воспроизводим по seed \\
\texttt{test\_alerting.py}     & уровни severity ladder; пропуск под порогом; дедупликация в окне; размер истории \\
\bottomrule
\end{tabularx}
\end{table}

\subsection{Соответствие реализации функциональным требованиям}

В \tabref{tab:traceability} приведено сводное соответствие
функциональных требований (FR-1…FR-10, сформулированных в
\secref{ch:theory}) реализованным артефактам и тестовому покрытию.

\begin{table}[H]
\centering
\caption{Соответствие реализации функциональным требованиям системы}
\label{tab:traceability}
\renewcommand{\arraystretch}{1.2}
\setlength{\tabcolsep}{4pt}
\footnotesize
\begin{tabularx}{\linewidth}{@{}p{2.6cm}Xp{2.8cm}@{}}
\toprule
\textbf{Требование} & \textbf{Артефакт реализации} & \textbf{Тест} \\
\midrule
FR-1: схема UNSW-NB15      & Pydantic-схема \texttt{schema.py}; валидация типов и диапазонов в \texttt{load.py}                        & \texttt{test\_configs} \\
FR-2: препроцессинг + окна & \texttt{Preprocessor}, \texttt{build\_windows}                                                            & \texttt{test\_preprocess}, \texttt{test\_windowing} \\
FR-3: модели ($\geq$ 14)   & 10 DL + 5 classical в едином registry \texttt{models/}                                                    & \texttt{test\_models\_smoke} \\
FR-4: порог + калибровка   & поиск порога по target-FPR; температурное масштабирование \texttt{TemperatureScaler}                      & \texttt{test\_metrics} \\
FR-5: формирование алертов & \texttt{AlertFormer} с severity ladder и temporal dedup                                                   & \texttt{test\_alerting} \\
FR-6: REST-API             & FastAPI-сервис \texttt{inference/service.py}, эндпоинты \texttt{/score, /info, /alerts}                   & ручное smoke \\
FR-7: визуальная панель    & Streamlit-дэшборд \texttt{ui/}                                                                            & ручное smoke \\
FR-8: ИИ-злоумышленник     & \texttt{FSMAgent}, \texttt{DriftInjector}, \texttt{AttackerRuntime}                                       & \texttt{test\_attacker} \\
FR-9: drift-monitoring     & функция \texttt{drift\_report} (PSI, KL, MMD)                                                             & \texttt{test\_metrics} \\
FR-10: seed + журнал       & \texttt{set\_seed}; JSON-логи в \texttt{experiments/}                                                     & \texttt{test\_utils} \\
\bottomrule
\end{tabularx}
\end{table}

\subsection*{Выводы по главе 3}

\begin{enumerate}
    \item Реализован программный прототип \texttt{diploma\_nids}
    объёмом порядка 4230 строк Python с восемью подпакетами и
    26 конфигурационными файлами.
    \item Реализованы 15 моделей-кандидатов (10 нейросетевых и 5
    классических) в едином registry с унифицированным интерфейсом.
    \item Реализован ИИ-злоумышленник в архитектуре
    \emph{templates + FSM + drift}, поддерживающий offline и online
    режимы и обеспечивающий воспроизводимую ground-truth.
    \item Реализованы подсистемы online-инференса (FastAPI),
    алертирования с severity ladder и дедупликацией, и визуальной
    панели (Streamlit), пригодные для эксплуатации в составе SOC.
    \item Покрытие тестами включает 9 модулей; критические инварианты
    pipeline проверены формально.
\end{enumerate}
